\chapter[Introdução]{Introdução}

Devido à convivência com o espaço de trabalho de uma lanchonete, observei algumas necessidades que poderiam ser automatizadas, já que o estabelecimento não estava operando com sua eficiência total, devido à falta de um sistema informatizado para controle de estoque, gerenciamento de vendas, gerenciamento de pedidos, geração de comandas, dentre outras funcionalidades.

A falta dessa ferramenta está ocasionando erros, atrasos e sobrecarga de trabalho para os funcionários, o que pode levar à insatisfação dos clientes e perda de vendas, ocasionando um ambiente de trabalho estressante e ineficiente, o que pode ser desagradável para donos, funcionários e clientes.

Com isso, vi a necessidade de desenvolver um sistema que pudesse otimizar o processo de trabalho e melhorar a eficácia, sem a necessidade de contratar mais funcionários, além de elevar o número de vendas na lanchonete, por meio do aperfeiçoamento do processo de trabalho.

Com o desenvolvimento do sistema, espera-se que venha trazer melhoria significativa para a lanchonete, otimizando o processo de trabalho, tornando um ambiente mais agradável, controlado e eficiente, além de buscar visar na satisfação dos clientes, funcionários e donos, o que pode levar a um aumento nas vendas e na fidelização dos clientes, além de melhorar a qualidade do serviço oferecido pela lanchonete, o que pode ser um diferencial competitivo no mercado.
